% Source: https://artofproblemsolving.com/community/c3h2416978_1200th_post_dedicated_to_bloggy
% Stable identity: aops-topic-2416978-q06
Write $12i$ in polar form: $12i=12\left(\cos90^\circ+i\sin90^\circ\right)$. Let $x=\rho\left(\cos\theta+i\sin\theta\right)$ with $\rho>0$. By De Moivre's theorem,
\[x^{12}=\rho^{12}\left(\cos12\theta+i\sin12\theta\right),\]
so we need $\rho^{12}=12$ and $12\theta=90^\circ+360^\circ k$ for an integer $k$. Thus $\rho=\sqrt[12]{12}$ and $\theta=7.5^\circ+30^\circ k$. The values $k=0,1,\dots,11$ give $12$ different numbers (further values of $k$ repeat them), and a polynomial equation of degree $12$ has at most $12$ roots, so these are all the solutions:
\[\boxed{x_k=\sqrt[12]{12}\left(\cos\left(7.5^\circ+30^\circ k\right)+i\sin\left(7.5^\circ+30^\circ k\right)\right),\quad k=0,1,\dots,11},\]
that is, $x_k=\sqrt[12]{12}\,e^{i\left(\dfrac{\pi}{24}+\dfrac{k\pi}{6}\right)}$. They are equally spaced, $30^\circ$ apart, on the circle of radius $\sqrt[12]{12}$.
